\section{统计公平到此为止}
\label{sec:17-Constrained-Guarantees/03-Fairness-Impossibility/03}

合规评审提了一个新问题，问法和前两节都不一样：``把这份申请里的群体标记改成另一个值，其余一切保持原样，模型的输出会不会变？''

工程师先做了最直接的实验：把 \(A\) 这一列从特征里删掉，重新训练。新模型的输出与旧模型几乎逐条一致，两者的相关系数是 \(0.98\)。有人据此宣布问题解决了——模型现在看不见群体属性。也有人据此宣布问题严重了——输出没变，说明群体信息本来就藏在别处。

两种解读都从同一个数字出发，而这个数字本身不能裁决。更麻烦的是，评审提的那个问题——改一个属性、其他不变、输出会怎样——根本不是前两节那些判据能回答的类型。

\subsection{删掉那一列为什么不起作用}

邮编、职业史、社交网络规模、消费类目，这些特征与 \(A\) 相关，模型可以从它们的组合里把 \(A\) 重建出来。删掉 \(A\) 只是删掉了那个显式的列，没有删掉信息。

这与\hyperref[sec:13-Machine-Learning/12-Causal-Inference/03]{``混杂、碰撞点与后门调整''}里讨论的是同一件事的两面。那里说的是：一个变量在数据表里的名字不决定它的角色，决定角色的是箭头怎么走。这里说的是：一个变量不在数据表里，不等于它的影响不在模型里——只要还有一组特征能张成含 \(A\) 的方向，模型在优化预测误差的过程中就会自动把它找出来，因为它对预测有用。\hyperref[sec:13-Machine-Learning/02-Linear-Models/02]{``正则化从 Ridge 到 Lasso''}里那条观察在这里同样适用：正则化压的是预测误差与系数规模，它认不出哪个方向对应敏感属性，也不会因为某个方向``不该被用''而绕开它。

真正的推论是：``模型不使用 \(A\)''这句话在统计层面几乎无法通过删列来实现，也无法通过检查特征清单来验证。要说清一个模型有没有使用群体信息，得先说清``使用''是什么意思，而这个词一旦被认真定义，就已经越出了联合分布的范围。

\subsection{所有统计判据都住在同一层}

回头看前两节的四个判据：人口平等、机会平等、等化几率、校准。它们的定义里只出现 \(P\)，参与运算的只有 \(A,S,\hat Y,Y\) 的联合分布。给定这个分布，四个判据的取值就完全确定了。

\hyperref[sec:13-Machine-Learning/12-Causal-Inference/01]{``观察、干预与反事实回答不同问题''}把问题分成三层：看见什么、做了什么之后会怎样、如果当初不同会怎样。四个判据全部住在第一层。而评审提的那个问题——``若这个人属于另一个群体，决定会不会变''——住在第三层，问的是同一个个体在另一个未发生的世界里的结果。

层与层之间不能靠算得更细来跨越。\hyperref[sec:13-Machine-Learning/12-Causal-Inference/02]{``结构因果模型把干预写成机制替换''}给出的理由是：干预与反事实是定义在机制上的，而联合分布只是机制的一个投影，多个不同的机制可以投影到同一个分布上。于是任何只读分布的判据，都无法区分那些分布相同、机制不同的世界。

\begin{center}
\begin{tikzpicture}[>=stealth,
  var/.style={draw,circle,minimum size=0.78cm,fill=blue!7},
  every node/.style={font=\small}]
  \begin{scope}
    \node[var] (a1) at (0,1.8) {\(A\)};
    \node[var] (r1) at (1.7,1.8) {\(R\)};
    \node[var] (d1) at (3.4,1.8) {\(D\)};
    \node[var] (y1) at (1.7,0.35) {\(Y\)};
    \draw[->] (a1) -- (r1);
    \draw[->] (r1) -- (d1);
    \draw[->] (r1) -- (y1);
    \node[align=center,font=\footnotesize] at (1.7,-0.75)
      {群体只通过真实风险起作用};
  \end{scope}
  \begin{scope}[xshift=6.4cm]
    \node[var] (a2) at (0,1.8) {\(A\)};
    \node[var] (r2) at (1.7,1.8) {\(R\)};
    \node[var] (d2) at (3.4,1.8) {\(D\)};
    \node[var] (y2) at (1.7,0.35) {\(Y\)};
    \draw[->] (a2) -- (r2);
    \draw[->] (r2) -- (d2);
    \draw[->] (r2) -- (y2);
    \draw[->,thick] (a2) to[bend left=42] (d2);
    \node[align=center,font=\footnotesize] at (1.7,-0.75)
      {群体还有一条直接进入决策的边};
  \end{scope}
  \node[anchor=west,align=left,font=\footnotesize,black!70] at (0,-1.6)
    {\(A\)：群体属性\quad \(R\)：真实风险\quad \(D\)：决策\quad \(Y\)：结果};
\end{tikzpicture}

\emph{两张图可以产生同一个观察分布，右图那条粗箭头的效应能被其余通路抵消，也能由一个未观测变量顶替。四个统计判据读的都是这个分布，因此对这条边说不出话。}
\end{center}

这张图要记住的是那条粗箭头的地位：它是否存在，改变了对系统的判断，却不改变数据。前两节那场争论至少还有共同的算术基础，双方各自算对了一个数；到了这一层，连能不能算都成了问题。

\subsection{反事实公平：把要求写在图上}

既然问题住在第三层，就用第三层的语言写要求。记 \(X\) 为其余特征，反事实公平要求：对每个个体，把 \(A\) 干预为 \(a\) 与干预为 \(a'\)，决策的分布相同，
\[
P\bigl(D_{A\leftarrow a}=d\given X=x,A=a\bigr)
=P\bigl(D_{A\leftarrow a'}=d\given X=x,A=a\bigr).
\]
下标记号是\hyperref[sec:13-Machine-Learning/12-Causal-Inference/04]{``潜在结果把处理效应定义在替代世界之间''}那一套：\(D_{A\leftarrow a}\) 是把 \(A\) 强行设为 \(a\) 之后这个个体的决策，两个世界只有一个个体、一次观测，另一个世界永远看不到。

在结构因果模型里，这个条件有一个可以执行的充分做法：让决策只依赖那些在图中不是 \(A\) 的后代的变量。这样一来，改变 \(A\) 不会沿任何有向路径传到 \(D\)，等式自动成立。这个做法之所以可执行，是因为它把一个反事实要求转成了一个关于图的结构要求。

为此要付出的是一张图。而图不能从数据里识别出来——这正是上面那两张图想说的事。所以要判断一个系统是否反事实公平，得先有人声明一张因果图，声明的依据是领域知识、制度知识与对数据生成过程的了解，不是这份数据本身。

这一步必须诚实地说出来：\textbf{公平性判断的争论没有消失，它被搬到了图上}。原先双方争的是该用哪个判据，现在双方争的是邮编到底算不算真实风险的合法度量、学历差异有多少来自可归因于系统的通路。这不是退步——把争论摆到明处、指着某个节点问``你确定这条边不存在吗''，比一句含糊的``模型是公平的''可检验得多。但它也不是解决，只是把不可判定的部分从暗处挪到了明处。

\hyperref[sec:13-Machine-Learning/12-Causal-Inference/06]{``可识别不等于可外推也不等于可决策''}里那条判断在这里再次生效：一个量在图上可识别，只说明给定图它能被算出来；图本身是否成立，是另一个层次的问题，而那个层次没有数据可以仲裁。

\subsection{这个单元留下什么}

三节走下来，可以固定几条判断。

四个统计判据各自对应一句朴素的正义直觉，而它们在对同一张 \(2\times2\) 表的不同分母做归一化，所以``模型公平吗''这个问句在被指定判据之前没有真值。这不是语言游戏，第一节那两份报告就是它的现实形态。

当两个群体的基率不同、且分类器会犯错时，校准、等真阳率、等假阳率不能同时成立，这一点由一个四行的代数恒等式给出。它的两个前提各挡住一种误读：基率相同则冲突消失，完美分类器则冲突消失。所以这条限制既不是模型的缺陷，也不能被更多数据推开——能选的只有让哪一条不成立。

而所有这些判据都住在观察层。人们真正想问的那类问题住在反事实层，跨过去需要一张因果图，图靠声明而不是靠数据。删掉敏感属性列不解决任何事，因为代理变量会把它重建出来。

还有一条属于整个系统而不属于模型的：这些判据检验的是模型对 \(Y\) 这份\emph{记录}的一致性，而记录本身经过了流程的过滤。系统上线后，它的判定又会改变谁被记录、以什么方式被记录，于是下一轮训练数据里的 \(Y\) 已经带着上一轮判定的痕迹——\hyperref[sec:18-Synthesis/01-System-Lifecycle/06]{``上线之后系统开始改变它所预测的世界''}讲的就是这种回路，\hyperref[sec:13-Machine-Learning/01-Learning-Problems-and-Evaluation/04]{``数据划分、泄漏与分布偏移''}给出了它在评估上的表现形式。任何静态的判据检查都看不见这条回路，因为它是跨时间的。

本部分其余单元里也会反复出现同一个形状：几件都想要的好事，被一条恒等式或一条界绑在一起，于是必须指定要哪一件。这一单元的特殊之处只在于，被绑住的那几件事都叫``公平''。
